\section{模型假设}

% 假设应有依据、可解释，并在后文说明其影响；不要罗列与模型无关的常识。
\begin{enumerate}
  \item 假设【数据记录或抽样过程】能够代表研究对象的主要特征；
  \item 假设【研究期内的关键关系】保持相对稳定；
  \item 假设【未观测因素】可通过误差项或敏感性分析进行评估。
\end{enumerate}

\section{符号说明}

% symboltable 为可跨页三线表；可选参数为表题，第三列填写单位或“--”。
\begin{symboltable}[主要符号说明]
  \label{tab:symbols}
  $x_{ij}$ & 第 $i$ 个样本在第 $j$ 个特征上的原始观测值 & -- \\
  $z_{ij}$ & 数据预处理后的特征值 & -- \\
  $\boldsymbol{\theta}$ & 待估计的模型参数向量 & -- \\
  $\hat{y}_i$ & 第 $i$ 个样本的预测值 & 题目给定单位 \\
  $L$ & 模型的损失函数或评价指标 & -- \\
\end{symboltable}
